\subsection{Equations}

Each equation can be expressed in either integral or differential form.
I'll present them in their differential form. We'll use 3D Cartesian
coordinates \((x,y,z)\) and time \(t\) to describe the functions.

\subsubsection{Faraday's Law of Induction}
Faraday's Law of Induction states that a time-varying magnetic function
induces an electric function. It is expressed as:
\begin{align*}
  -\nabla \times \mathbf{E(x,y,z,t)} = \frac{\partial \mathbf{B(x,y,z,t)}}{\partial t} + \mathbf{J_m(x,y,z,t)} \tag{1}
\end{align*}
where
\begin{itemize}
  \item \( \nabla \):
    \begin{align*}
      \nabla = \left( \frac{\partial}{\partial x}, \frac{\partial}{\partial y}, \frac{\partial}{\partial z} \right)
    \end{align*}
  \item \( \mathbf{E(x,y,z,t)} \): electric field function. \( \frac{V}{m} \)
  \item \( \mathbf{B(x,y,z,t)} \): magnetic flux density function.
    \( \frac{Wb}{m^2} \)
  \item \( \mathbf{J_m(x,y,z,t)} \): magnetic current density function.
    \( \frac{V}{m^2} \), the same unit as \(\frac{\partial \mathbf{B}}{\partial t}\), which is
    \(\frac{Wb}{m^2 s}\). In nature, \( \mathbf{J_m(x,y,z,t)} = 0 \) since
    magnetic monopoles that is, charges haven't been observed.
\end{itemize}
Let's have some fun with this equation.
\begin{align*}
  & \text{[Divergence of both sides]} \\
  -\nabla \cdot (\nabla \times \mathbf{E(x,y,z,t)}) &= \nabla \cdot \left(\frac{\partial \mathbf{B(x,y,z,t)}}{\partial t} + \mathbf{J_m(x,y,z,t)}\right) \\
  -\nabla \cdot (\nabla \times \mathbf{E(x,y,z,t)}) &= \nabla \cdot \frac{\partial \mathbf{B(x,y,z,t)}}{\partial t} + \nabla \cdot \mathbf{J_m(x,y,z,t)} \\
  -\nabla \cdot (\nabla \times \mathbf{E(x,y,z,t)}) &= \frac{\partial}{\partial t} (\nabla \cdot \mathbf{B(x,y,z,t)}) + \nabla \cdot \mathbf{J_m(x,y,z,t)} \\
  0 &= \frac{\partial}{\partial t} (\nabla \cdot \mathbf{B(x,y,z,t)}) + \nabla \cdot \mathbf{J_m(x,y,z,t)} \\
  - \frac{\partial}{\partial t} (\nabla \cdot \mathbf{B(x,y,z,t)}) &= \nabla \cdot \mathbf{J_m(x,y,z,t)} \\
  & \text{[Integrating both sides over time]} \\
  \int_{t_0}^{t_1} - \frac{\partial}{\partial t} (\nabla \cdot \mathbf{B(x,y,z,t)}) \, dt &= \int_{t_0}^{t_1} \nabla \cdot \mathbf{J_m(x,y,z,t)} \, dt \\
  - \left[\nabla \cdot \mathbf{B(x,y,z,t)}\right]_{t_0}^{t_1} &= \int_{t_0}^{t_1} \nabla \cdot \mathbf{J_m(x,y,z,t)} \, dt \\
  - \left[\nabla \cdot \mathbf{B(x,y,z,t_1)} - \nabla \cdot \mathbf{B(x,y,z,t_0)}\right] &= \int_{t_0}^{t_1} \nabla \cdot \mathbf{J_m(x,y,z,t)} \, dt \\
  - \nabla \cdot \mathbf{B(x,y,z,t_1)} + \nabla \cdot \mathbf{B(x,y,z,t_0)} &=  \int_{t_0}^{t_1} \nabla \cdot \mathbf{J_m(x,y,z,t)} \, dt \\
  - \nabla \cdot \mathbf{B(x,y,z,t_1)} &= - \nabla \cdot \mathbf{B(x,y,z,t_0)} + \int_{t_0}^{t_1} \nabla \cdot \mathbf{J_m(x,y,z,t)} \, dt \\
  - \nabla \cdot \mathbf{B(x,y,z,t_1)} &= - \left[\nabla \cdot \mathbf{B(x,y,z,t_0)} - \int_{t_0}^{t_1} \nabla \cdot \mathbf{J_m(x,y,z,t)} \, dt\right] \\
  \nabla \cdot \mathbf{B(x,y,z,t_1)} &= \nabla \cdot \mathbf{B(x,y,z,t_0)} - \int_{t_0}^{t_1} \nabla \cdot \mathbf{J_m(x,y,z,t)} \, dt \\
  & \text{[Since \( \mathbf{J_m(x,y,z,t)} = 0 \) in nature]} \\
  \nabla \cdot \mathbf{B(x,y,z,t_1)} &= \nabla \cdot \mathbf{B(x,y,z,t_0)} - \int_{t_0}^{t_1} 0 \, dt \\
  \nabla \cdot \mathbf{B(x,y,z,t_1)} &= \nabla \cdot \mathbf{B(x,y,z,t_0)}
\end{align*}
\textbf{This proves that the divergence of the magnetic flux density function
is constant over time}; it does not fix its value. Gauss's Law for Magnetism
is a separate law (no magnetic charge has ever been observed), and Faraday's
law only guarantees that once it holds at some \(t_0\) it holds for all
time. With \(\nabla \cdot \mathbf{B} = 0\) at \(t_0\), the result above keeps it zero:
\begin{align*}
  \nabla \cdot \mathbf{B(x,y,z,t)} = 0 \tag{2}
\end{align*}

\subsubsection{Ampere's Law}

Ampere's Law states that a time-varying electric function induces a magnetic
function. It is expressed as:
\begin{align*}
  \nabla \times \mathbf{H(x,y,z,t)} = \frac{\partial \mathbf{D(x,y,z,t)}}{\partial t} + \mathbf{J_e(x,y,z,t)} \tag{3}
\end{align*}
where
\begin{itemize}
  \item \( \mathbf{H(x,y,z,t)} \): magnetic function. \( \frac{A}{m} \)
  \item \( \mathbf{D(x,y,z,t)} \): electric flux density function.
    \( \frac{C}{m^2} \)
  \item \( \mathbf{J_e(x,y,z,t)} \): electric current density function.
    \( \frac{A}{m^2} \). In nature, \( \mathbf{J_e(x,y,z,t)} \neq 0 \)
    since electric charges exist.
    \begin{align*}
      \mathbf{J_e(x,y,z,t)} = Q_V(x,y,z,t) \mathbf{v(x,y,z,t)}
    \end{align*}
    with \(Q_V\) the volume charge density (\(\frac{C}{m^3}\)) and \(\mathbf{v}\) the
    charge velocity (\(\frac{m}{s}\)).
\end{itemize}
Let's have some fun with this equation.
\begin{align*}
  & \text{[Divergence of both sides]} \\
  \nabla \cdot (\nabla \times \mathbf{H(x,y,z,t)}) &= \nabla \cdot \left(\frac{\partial \mathbf{D(x,y,z,t)}}{\partial t} + \mathbf{J_e(x,y,z,t)}\right) \\
  \nabla \cdot (\nabla \times \mathbf{H(x,y,z,t)}) &= \nabla \cdot \frac{\partial \mathbf{D(x,y,z,t)}}{\partial t} + \nabla \cdot \mathbf{J_e(x,y,z,t)} \\
  \nabla \cdot (\nabla \times \mathbf{H(x,y,z,t)}) &= \frac{\partial}{\partial t} (\nabla \cdot \mathbf{D(x,y,z,t)}) + \nabla \cdot \mathbf{J_e(x,y,z,t)} \\
  0 &= \frac{\partial}{\partial t} (\nabla \cdot \mathbf{D(x,y,z,t)}) + \nabla \cdot \mathbf{J_e(x,y,z,t)} \\
  - \frac{\partial}{\partial t} (\nabla \cdot \mathbf{D(x,y,z,t)}) &= \nabla \cdot \mathbf{J_e(x,y,z,t)} \\
  & \text{[Integrating both sides over time]} \\
  \int_{t_0}^{t_1} - \frac{\partial}{\partial t} (\nabla \cdot \mathbf{D(x,y,z,t)}) \, dt &= \int_{t_0}^{t_1} \nabla \cdot \mathbf{J_e(x,y,z,t)} \, dt \\
  - \left[\nabla \cdot \mathbf{D(x,y,z,t)}\right]_{t_0}^{t_1} &= \int_{t_0}^{t_1} \nabla \cdot \mathbf{J_e(x,y,z,t)} \, dt \\
  - \left[\nabla \cdot \mathbf{D(x,y,z,t_1)} - \nabla \cdot \mathbf{D(x,y,z,t_0)}\right] &= \int_{t_0}^{t_1} \nabla \cdot \mathbf{J_e(x,y,z,t)} \, dt \\
  - \nabla \cdot \mathbf{D(x,y,z,t_1)} + \nabla \cdot \mathbf{D(x,y,z,t_0)} &=  \int_{t_0}^{t_1} \nabla \cdot \mathbf{J_e(x,y,z,t)} \, dt \\
  - \nabla \cdot \mathbf{D(x,y,z,t_1)} &= - \nabla \cdot \mathbf{D(x,y,z,t_0)} + \int_{t_0}^{t_1} \nabla \cdot \mathbf{J_e(x,y,z,t)} \, dt \\
  - \nabla \cdot \mathbf{D(x,y,z,t_1)} &= - \left[\nabla \cdot \mathbf{D(x,y,z,t_0)} - \int_{t_0}^{t_1} \nabla \cdot \mathbf{J_e(x,y,z,t)} \, dt \right] \\
  \nabla \cdot \mathbf{D(x,y,z,t_1)} &= \nabla \cdot \mathbf{D(x,y,z,t_0)} - \int_{t_0}^{t_1} \nabla \cdot \mathbf{J_e(x,y,z,t)} \, dt
\end{align*}
\textbf{This proves that the divergence of the electric flux density function
changes over time}. To find anything meaningful we must find more information
about the divergence of the electric current density function. This is given
to us by the \textbf{charge continuity equation}:
\begin{align*}
  & \text{[Total Charge in a Volume]} \\
  Q(t) &= \int_V Q_V(x,y,z,t) \, dV \\
  & \text{[Rate of Change of Total Charge]} \\
  & \text{[In a Volume]} \\
  \frac{dQ(t)}{dt} &= \int_V \frac{\partial Q_V(x,y,z,t)}{\partial t} \, dV \\
  & \text{[Out of a Volume]} \\
  \frac{dQ(t)}{dt} &= - \oint_S \mathbf{J_e(x,y,z,t)} \cdot d\mathbf{S} \\
  & \text{[Continuity; In = Out]} \\
  - \oint_S \mathbf{J_e(x,y,z,t)} \cdot d\mathbf{S} &= \int_V \frac{\partial Q_V(x,y,z,t)}{\partial t} \, dV \\
  & \text{[Divergence theorem: } \oint_S \mathbf{J_e} \cdot d\mathbf{S} = \int_V \nabla \cdot \mathbf{J_e} \, dV\text{]} \\
  - \int_V \nabla \cdot \mathbf{J_e(x,y,z,t)} \, dV &= \int_V \frac{\partial Q_V(x,y,z,t)}{\partial t} \, dV \\
  \int_V \left(- \nabla \cdot \mathbf{J_e(x,y,z,t)} - \frac{\partial Q_V(x,y,z,t)}{\partial t}\right) dV &= 0 \\
  & \text{[True for every volume \(V\), however small,} \\
  & \text{\ so the integrand is zero]} \\
  - \nabla \cdot \mathbf{J_e(x,y,z,t)} &= \frac{\partial Q_V(x,y,z,t)}{\partial t} \tag{4}
\end{align*}
Substituting (4) into the time integrated result above, we get
\begin{align*}
  \nabla \cdot \mathbf{D(x,y,z,t_1)} &= \nabla \cdot \mathbf{D(x,y,z,t_0)} - \int_{t_0}^{t_1} \nabla \cdot \mathbf{J_e(x,y,z,t)} \, dt \\
  \nabla \cdot \mathbf{D(x,y,z,t_1)} &= \nabla \cdot \mathbf{D(x,y,z,t_0)} - \int_{t_0}^{t_1} - \frac{\partial Q_V(x,y,z,t)}{\partial t} \, dt \\
  \nabla \cdot \mathbf{D(x,y,z,t_1)} &= \nabla \cdot \mathbf{D(x,y,z,t_0)} + \int_{t_0}^{t_1} \frac{\partial Q_V(x,y,z,t)}{\partial t} \, dt \\
  \nabla \cdot \mathbf{D(x,y,z,t_1)} &= \nabla \cdot \mathbf{D(x,y,z,t_0)} + [Q_V(x,y,z,t_1) - Q_V(x,y,z,t_0)] \\
  \nabla \cdot \mathbf{D(x,y,z,t_1)} &= \nabla \cdot \mathbf{D(x,y,z,t_0)} + Q_V(x,y,z,t_1) - Q_V(x,y,z,t_0)
\end{align*}
Rearranged, \(\nabla \cdot \mathbf{D}(t_1) - Q_V(t_1) = \nabla \cdot \mathbf{D}(t_0) - Q_V(t_0)\): the
difference \(\nabla \cdot \mathbf{D} - Q_V\) never changes. Ampere's law with charge
continuity therefore does not produce Gauss's Law of Electricity, it only
preserves it. Gauss's law is the separate statement that this difference is
zero; once it holds at any \(t_0\) it holds for all time:
\begin{align*}
  \nabla \cdot \mathbf{D(x,y,z,t)} &= Q_V(x,y,z,t) \tag{3a}
\end{align*}

\subsubsection{Statics \& Quasi-Statics Approximation}

Faraday's and Ampere's Law (alongside Gauss Conditions) fully define any electromagnetic system
at all times and space. Both equations are \textbf{first order partial differential
equations (PDEs)} which \textbf{are coupled with each other}. Sadly, for solving the
full system is mathematically very expensive. One therefore wonders, \textbf{Is it
possible to simplify and dumb down?}


\textbf{Yes, but only under conditions}. This is where Statics and Quasi-Statics
come into play. We need to make some assumptions.
\begin{enumerate}
  \item \textbf{Statics}: We assume that all functions are frozen in time. This is the
    strongest assumption we can make and it removes \(t\), leaving \textbf{PDEs in
    space only} (\(x, y, z\); e.g.\ \(\nabla \times \mathbf{E} = 0\) with \(\nabla \cdot \mathbf{D} = Q_V\)). This is the DC regime where all functions
    are time-independent. To achieve this, the rate of change of all functions must
    be zero:
    \begin{align*}
      \frac{\partial \mathbf{B(x,y,z,t)}}{\partial t} = 0 \\
      \frac{\partial \mathbf{D(x,y,z,t)}}{\partial t} = 0 \\
      \frac{\partial \mathbf{E(x,y,z,t)}}{\partial t} = 0 \\
      \frac{\partial \mathbf{H(x,y,z,t)}}{\partial t} = 0 \\
      \frac{\partial \mathbf{J_e(x,y,z,t)}}{\partial t} = 0 \\
      \frac{\partial \mathbf{J_m(x,y,z,t)}}{\partial t} = 0
    \end{align*}
  \item \textbf{Quasi-Statics}: We assume that the rates of change of the flux
    functions [\(\mathbf{B}\) \& \(\mathbf{D}\)] are negligible next to the other terms of
    Faraday's and Ampere's laws, while we allow the other functions
    [\(\mathbf{E}\), \(\mathbf{H}\), \(\mathbf{J}_e\), and \(\mathbf{J}_m\)] to change with time. The
    flux functions cannot be exactly constant (with \(\mathbf{D} = \mathbf{K}_E \mathbf{E}\) a frozen
    \(\mathbf{D}\) would freeze \(\mathbf{E}\) too), so this is an approximation: it holds when the
    region is much smaller than the wavelength, \(2\pi\frac{l}{\text{wavelength}} \ll 1\) for a
    region of size \(l\) (e.g.\ \(100\,MHz\): wavelength \(= \frac{3 \times 10^{8}\,m/s}{10^{8}\,Hz} = 3\,m\), so a
    \(3\,cm\) board gives \(2\pi \cdot \frac{0.03}{3} = 0.063 \ll 1\)). This is a weaker assumption and,
    at each instant, it leaves PDEs in space only. This is the quasi-static
    regime. In a circuit it is applied \emph{between} the elements: all of the
    \(\frac{\partial \mathbf{B}}{\partial t}\) is kept inside inductors (\(V = L\frac{dI}{dt}\)) and all of the
    \(\frac{\partial \mathbf{D}}{\partial t}\) inside capacitors (\(I = C\frac{dV}{dt}\)). Everywhere else:
    \begin{align*}
      \frac{\partial \mathbf{B(x,y,z,t)}}{\partial t} \approx 0 \\
      \frac{\partial \mathbf{D(x,y,z,t)}}{\partial t} \approx 0
    \end{align*}
\end{enumerate}
In either case, we can simplify Faraday's and Ampere's Law to:
\begin{align*}
    -\nabla \times \mathbf{E(x,y,z,t)} &= \frac{\partial \mathbf{B(x,y,z,t)}}{\partial t} + \mathbf{J_m(x,y,z,t)} \hspace{1cm} \text{[Faraday's Law]} \\
                                       &\approx 0 + \mathbf{J_m(x,y,z,t)} \hspace{1cm} \left[\frac{\partial \mathbf{B(x,y,z,t)}}{\partial t} \approx 0\right] \\
                                       &= \mathbf{J_m(x,y,z,t)} \\
     &= 0  \hspace{1cm} \text{[since \( \mathbf{J_m(x,y,z,t)} = 0 \)]} \tag{5} \\\\
    \nabla \times \mathbf{H(x,y,z,t)} &= \frac{\partial \mathbf{D(x,y,z,t)}}{\partial t} + \mathbf{J_e(x,y,z,t)} \hspace{1cm} \text{[Ampere's Law]} \\
                                      &\approx 0 + \mathbf{J_e(x,y,z,t)} \hspace{1cm} \left[\frac{\partial \mathbf{D(x,y,z,t)}}{\partial t} \approx 0\right] \\
                                      &= \mathbf{J_e(x,y,z,t)} \tag{6}
\end{align*}

Equations (5) and (6) are \textbf{instantaneous spatial statements}. They hold at every
value of \(t\) individually but say nothing about how \(\mathbf{E}\) or \(\mathbf{H}\)
behave \emph{between} t's. A function can satisfy \(\nabla\times\mathbf{E}=0\) at
every instant and still change with time. For example \(\mathbf{E} = \sin(\omega t)\,
\nabla\phi(x,y,z)\) has zero curl for all \(t\), yet clearly changes with time. So (5)
and (6) do not, by themselves, let us drop \(t\) from \(\mathbf{E}\), \(\mathbf{H}\),
or \(\mathbf{J}_e\) (\& \(\mathbf{J}_m\)).

\begin{enumerate}
  \item \textbf{Equation 5}: At every instant (t), the electric function  has zero
    curl = it is irrotational = it is conservative = the potential difference between
    two points does not depend on the path taken between them = total work done by
    the electric function on a charge moving from point a to point b and back to
    point a is zero, \emph{at that instant}.
    \begin{align*}
      & \text{[At every instant \(t\)]} \\
      \nabla \times \mathbf{E(x,y,z,t)} &= 0 \\
      & \text{[A curl free field on a region without holes is a gradient; define } V \text{ by]} \\
      & \text{[path independent: zero curl, Stokes]} \\
      V(\mathbf{r},t) &= -\int_{\mathbf{r}_0}^{\mathbf{r}} \mathbf{E} \cdot d\mathbf{l} \\
      & \text{[check: } \nabla \times (-\nabla V) = 0 \text{, curl of a gradient]} \\
      \mathbf{E(x,y,z,t)} &= -\nabla V(x,y,z,t) \\
      & \text{[Line Integral of Electric Function: A to B to A]} \\
      \oint_a^b \mathbf{E(x,y,z,t)} \cdot d\mathbf{l} &= \oint_a^b -\nabla V(x,y,z,t) \cdot d\mathbf{l} \\
                                                &= \int_a^b -\nabla V(x,y,z,t) \cdot d\mathbf{l} + \int_b^a -\nabla V(x,y,z,t) \cdot d\mathbf{l} \\
                                                &= -\left[ V(b,t) - V(a,t) \right] - \left[ V(a,t) - V(b,t) \right] \\
                                                &= 0
    \end{align*}
    This is the origin of Kirchhoff's Voltage Law (KVL): the total potential difference
    around any closed loop is zero, with each inductor's \(L\frac{dI}{dt}\) counted as
    the voltage across its terminals (the loop runs outside its flux). Notice this did not require \(\mathbf{E}\) or \(V\)
    to be time-independent, only that \(\nabla\times\mathbf{E}=0\) holds instant by
    instant. This is precisely why KVL remains valid in the quasi-static (slowly
    time-varying, circuit size much smaller than the wavelength) regime, not just in true DC.
  \item \textbf{Equation 6}: Taking the divergence of both sides of
    equation (6), at each instant \(t\):
    \begin{align*}
      \nabla \cdot (\nabla \times \mathbf{H(x,y,z,t)}) &= \nabla \cdot \mathbf{J_e(x,y,z,t)} \\
      0 &= \nabla \cdot \mathbf{J_e(x,y,z,t)} \hspace{1cm} \text{[since divergence of a curl is always zero]} \tag{7}
    \end{align*}

    \begin{itemize}
      \item \textbf{Charge Density Is Time-Independent (between capacitors)}:
        Substituting (4) into (7):
        \begin{align*}
          0 &= -\frac{\partial Q_V(x,y,z,t)}{\partial t} \tag{8}
        \end{align*}
        Equation (8) \textbf{shows the total charge enclosed by that
        surface genuinely does not change with time.} Unlike
        \(\mathbf{E}\), \(\mathbf{H}\), and \(\mathbf{J}_e\) above,
        this is a genuine, non-instantaneous result: it says
        \(Q_V\)'s rate of change is zero \emph{for all} \(t\), not
        just at a single instant, everywhere \(\frac{\partial \mathbf{D}}{\partial t}\) was neglected
        (wires and nodes). On capacitor plates, where \(\frac{\partial \mathbf{D}}{\partial t}\) was kept,
        charge does change, which is the \(C\frac{dV}{dt}\) current. Between capacitors
        this is the one function we can honestly write without \(t\):
        \begin{align*}
          Q_V(x,y,z,t) = Q_V(x,y,z) \tag{8a}
        \end{align*}
      \item \textbf{Zero Total Current Out of a Volume}: Integrating
        both sides of equation (7) over a volume \(V\), at a fixed
        instant \(t\):
        \begin{align*}
          \int_V \nabla \cdot \mathbf{J_e(x,y,z,t)} \, dV &= \int_V 0 \, dV \\
          \int_V \nabla \cdot \mathbf{J_e(x,y,z,t)} \, dV &= 0 \\
          \oint_S \mathbf{J_e(x,y,z,t)} \cdot d\mathbf{S} &= 0 \hspace{1cm} \text{[by Divergence Theorem]} \tag{9}
        \end{align*}
        Equation (9) \textbf{shows that the total current flowing out of a
        closed surface is zero at every instant \(t\)} without requiring
        \(\mathbf{J}_e\) itself to be constant in time.
    \end{itemize}
    If current can not flow out of a closed surface (9) and the total charge
    enclosed by that surface does not change with time (8a), then the total
    current flowing into that surface must equal the total current flowing
    out of that surface, at every instant (as current density is allowed to
    change with time). This is the origin of Kirchhoff's Current Law (KCL),
    with each capacitor's \(C\frac{dV}{dt}\) counted as the current in its leads.
\end{enumerate}
